\chapter{ผลการดำเนินงานรายโครงการหลัก}

บทนี้นำเสนอผลการดำเนินงานของโครงการหลักทั้งสามตามกรอบ Logic Model
คือแยกเป็น \textbf{ผลผลิต} (สิ่งที่ผลิตหรือส่งมอบได้) \textbf{ผลลัพธ์} (การเปลี่ยนแปลงที่เกิดกับผู้รับประโยชน์)
และ \textbf{ผลกระทบ} (ผลระยะยาวต่อพื้นที่และมหาวิทยาลัย)

\begin{infobox}
ผลผลิตและผลลัพธ์ที่ระบุในบทนี้อ้างอิงจาก\textbf{ตัวชี้วัดที่แต่ละโครงการรับมาดำเนินการ}
ตามไฟล์จัดสรรแผนงานและตัวชี้วัดฉบับ 20 กรกฎาคม 2569
ส่วนผลกระทบเป็นการวิเคราะห์เชิงคุณภาพจากลักษณะงานและพื้นที่ดำเนินการ
\end{infobox}

\section{ง8-1 ผลักดันเทคโนโลยี นวัตกรรมสู่ชุมชน ตามเป้าหมายการพัฒนาอย่างยั่งยืน}
16~โครงการย่อย $\cdot$ 6~หน่วยงานร่วม $\cdot$ งบจัดสรร 1,988,000~บาท $\cdot$ เบิกจ่าย 1,064,116~บาท (54\%)

\begin{table}[H]
\centering
\caption{ผลผลิต ผลลัพธ์ และผลกระทบของโครงการหลัก ง8-1}
\small
\begin{tabular}{L{2.2cm}L{11.5cm}}
\toprule
\rowcolor{headerblue}
\textbf{ระดับผล} & \textbf{สาระสำคัญ}\\
\midrule
\textbf{ผลผลิต} & เครื่องต้นแบบและเทคโนโลยีการผลิตและแปรรูปบนพื้นที่สูง ; ทรัพย์สินทางปัญญาที่จะยื่นจดทะเบียน 8~ผลงาน ; องค์ความรู้ส่งมอบพื้นที่โครงการหลวง 26~รายการ\\
\rowcolor{rowblue}
\textbf{ผลลัพธ์} & ศูนย์พัฒนาโครงการหลวงและชุมชนใช้เทคโนโลยีเพื่อลดต้นทุนและเพิ่มคุณภาพผลผลิต ; บุคลากร 98~คน นำเทคโนโลยีลงพื้นที่จริง\\
\textbf{ผลกระทบ} & ยกระดับรายได้และขีดความสามารถของเกษตรกรบนพื้นที่สูง ; สร้างฐานทรัพย์สินทางปัญญาให้มหาวิทยาลัย\\
\bottomrule
\end{tabular}
\end{table}

ตัวอย่างงานในกลุ่มนี้ ได้แก่ โครงการวิจัย Plant profile ของผักในโรงเรือนอัจฉริยะ
โครงการพัฒนากระบวนการสกัดน้ำมันอะโวคาโดด้วยเทคโนโลยีไฮดรอลิก
โครงการพัฒนาเทคโนโลยีเครื่องล้างเมล็ดงาดำแบบกึ่งอัตโนมัติ
และโครงการพัฒนาตู้อบลมร้อนดอกเก๊กฮวยระบบปั๊มความร้อน

\section{ง8-2 ขับเคลื่อนกลไกการพัฒนาองค์ความรู้เพื่อยกระดับคุณภาพชีวิต}
14~โครงการย่อย $\cdot$ 8~หน่วยงานร่วม $\cdot$ งบจัดสรร 1,999,010~บาท $\cdot$ เบิกจ่าย 1,360,534~บาท
(68\% --- สูงสุดในแผนงาน)

\begin{table}[H]
\centering
\caption{ผลผลิต ผลลัพธ์ และผลกระทบของโครงการหลัก ง8-2}
\small
\begin{tabular}{L{2.2cm}L{11.5cm}}
\toprule
\rowcolor{headerblue}
\textbf{ระดับผล} & \textbf{สาระสำคัญ}\\
\midrule
\textbf{ผลผลิต} & แหล่งเรียนรู้ตลอดชีวิต 4~แหล่ง ; พื้นที่เรียนรู้ศาสตร์พระราชาและเกษตรทฤษฎีใหม่ ; องค์ความรู้ส่งมอบพื้นที่ 13~รายการ\\
\rowcolor{rowblue}
\textbf{ผลลัพธ์} & ชุมชนเข้าถึงองค์ความรู้ผ่านแหล่งเรียนรู้ถาวร ; บุคลากร 58~คน ถ่ายทอดความรู้ในพื้นที่ ; ศูนย์พัฒนาพันธุ์พืชจักรพันธ์เพ็ญศิริและงานรับเสด็จได้รับการสนับสนุนต่อเนื่อง\\
\textbf{ผลกระทบ} & เกิดกลไกการเรียนรู้ที่ยั่งยืนไม่ผูกกับรอบปีงบประมาณ ; คุณภาพชีวิตชุมชนดีขึ้นตามหลักปรัชญาเศรษฐกิจพอเพียง\\
\bottomrule
\end{tabular}
\end{table}

ตัวอย่างงานในกลุ่มนี้ ได้แก่ โครงการพัฒนาพื้นที่เรียนรู้ศาสตร์พระราชา
โครงการสนับสนุนการดำเนินงานศูนย์พัฒนาพันธุ์พืชจักรพันธ์เพ็ญศิริ
โครงการพัฒนาแหล่งเรียนรู้เกษตรอันเนื่องมาจากพระราชดำริ
และโครงการถ่ายทอดองค์ความรู้การจัดการขยะแบบยั่งยืน

\section{ง8-3 พัฒนากำลังคน สร้างอาชีพ ลดความเหลื่อมล้ำ บนพื้นที่สูง}
31~โครงการย่อย $\cdot$ 9~หน่วยงานร่วม $\cdot$ งบจัดสรร 3,999,573~บาท (ร้อยละ~50 ของงบแผนงาน)
$\cdot$ เบิกจ่าย 2,035,682~บาท (51\%)

\begin{table}[H]
\centering
\caption{ผลผลิต ผลลัพธ์ และผลกระทบของโครงการหลัก ง8-3}
\small
\begin{tabular}{L{2.2cm}L{11.5cm}}
\toprule
\rowcolor{headerblue}
\textbf{ระดับผล} & \textbf{สาระสำคัญ}\\
\midrule
\textbf{ผลผลิต} & หลักสูตรและการอบรมทักษะอาชีพ Re-skill Up-skill New-skill ; เครือข่ายความร่วมมือ 7~เครือข่าย ; องค์ความรู้ส่งมอบพื้นที่ 23~รายการ\\
\rowcolor{rowblue}
\textbf{ผลลัพธ์} & เยาวชน สามเณร เจ้าหน้าที่ และชาวบ้านมีทักษะอาชีพใหม่ ; บุคลากร 132~คน ลงพื้นที่พัฒนากำลังคน ; วิสาหกิจชุมชนยกระดับมาตรฐานผลิตภัณฑ์\\
\textbf{ผลกระทบ} & ลดความเหลื่อมล้ำทางโอกาสและรายได้บนพื้นที่สูง ; สร้างกำลังคนที่อยู่กับชุมชนได้ในระยะยาว\\
\bottomrule
\end{tabular}
\end{table}

ตัวอย่างงานในกลุ่มนี้ ได้แก่ โครงการพัฒนาทักษะอาชีพสำหรับสามเณรในโรงเรียนพระปริยัติธรรม
โครงการยกระดับมาตรฐานผลิตภัณฑ์ผ้าใยกัญชงวิสาหกิจชุมชนดาวม่าง
โครงการยกระดับห่วงโซ่คุณค่าห้อมพื้นถิ่น
และโครงการพัฒนาทักษะอาชีพด้านพลังงานสะอาด

\section{สรุปเปรียบเทียบทั้งสามโครงการหลัก}
เมื่อพิจารณาร่วมกัน ง8-1 ทำหน้าที่เป็นแกน\textbf{ด้านเทคโนโลยีและทรัพย์สินทางปัญญา}
ง8-2 เป็นแกน\textbf{ด้านองค์ความรู้และความยั่งยืนของกลไกการเรียนรู้} และ ง8-3 เป็นแกน
\textbf{ด้านคนและอาชีพ} ซึ่งใหญ่ที่สุดทั้งจำนวนโครงการและวงเงิน
การแบ่งบทบาทนี้สอดคล้องกับการกระจายตัวชี้วัดที่นำเสนอในบทที่ 4
